%% ma: L11-B06-0012
%% lop: 11
%% chuong: 2
%% bai: 6
%% dang: 11.6.10
%% loai: TLN
%% muc: VD
%% dap_an: "2380"
%% nguon: "de-cuong-giua-ky-1-toan-11-nam-2026-2027-truong-thpt-hoang-van-thu-ha-noi-2-8.pdf (M08, Đề cương giữa kì 1 Toán 11 THPT Hoàng Văn Thụ, Hà Nội), Phần 3 Câu 7"
%% kien_thuc: "Bài 6 (tổng n số hạng đầu của cấp số cộng)"
%% trang_thai: cho-duyet
%% ly_do: "Dạng mới đề xuất 11.6.10"
%% ghi_chu: "Hiểu viên gạch dài 0,2 m theo hàng và cao 0,08 m theo chiều cao ngôi nhà (34 hàng × 0,08 m = 2,72 m). Kích thước thứ ba (1 m) không dùng đến; giữ nguyên như gốc. Tài liệu gốc không có hình."
\begin{ex}
Một ngôi nhà hình kim tự tháp (có gạch nâu ốp bên ngoài) được bao quanh bởi rất nhiều cây cối và là nơi tuyệt vời để nghỉ mát mùa hè; ngôi nhà có chiều dài, chiều rộng là $6{,}8$ m, chiều cao là $2{,}72$ m. Khi xây dựng ngôi nhà, người chủ đã tính toán số viên gạch nâu hình hộp chữ nhật cần ốp tường; biết hàng trên ít hơn hàng dưới $1$ viên, hàng trên cùng là $1$ viên, kích thước viên gạch nâu hình hộp chữ nhật là $0{,}2\text{ m}\times0{,}08\text{ m}\times1\text{ m}$. Hãy dự tính số viên gạch nâu ốp tường cả bốn mặt của ngôi nhà.
\shortans{2380}
\loigiai{Hàng dưới cùng mỗi mặt có $\dfrac{6{,}8}{0{,}2}=34$ viên; số hàng là $\dfrac{2{,}72}{0{,}08}=34$, khớp với các hàng $1,2,\dots,34$ viên.

Mỗi mặt có $1+2+\dots+34=\dfrac{34\cdot35}{2}=595$ viên. Bốn mặt: $4\cdot595=2380$ viên.}
\end{ex}
